% Options for packages loaded elsewhere
\PassOptionsToPackage{unicode}{hyperref}
\PassOptionsToPackage{hyphens}{url}
\PassOptionsToPackage{dvipsnames,svgnames,x11names}{xcolor}
\documentclass[
  11pt,
]{article}
\usepackage{xcolor}
\usepackage[left=3cm,right=3in,top=2cm,bottom=2cm]{geometry}
\usepackage{amsmath,amssymb}
\setcounter{secnumdepth}{5}
\usepackage{iftex}
\ifPDFTeX
  \usepackage[T1]{fontenc}
  \usepackage[utf8]{inputenc}
  \usepackage{textcomp} % provide euro and other symbols
\else % if luatex or xetex
  \usepackage{unicode-math} % this also loads fontspec
  \defaultfontfeatures{Scale=MatchLowercase}
  \defaultfontfeatures[\rmfamily]{Ligatures=TeX,Scale=1}
\fi
\usepackage{lmodern}
\ifPDFTeX\else
  % xetex/luatex font selection
\fi
% Use upquote if available, for straight quotes in verbatim environments
\IfFileExists{upquote.sty}{\usepackage{upquote}}{}
\IfFileExists{microtype.sty}{% use microtype if available
  \usepackage[]{microtype}
  \UseMicrotypeSet[protrusion]{basicmath} % disable protrusion for tt fonts
}{}
\makeatletter
\@ifundefined{KOMAClassName}{% if non-KOMA class
  \IfFileExists{parskip.sty}{%
    \usepackage{parskip}
  }{% else
    \setlength{\parindent}{0pt}
    \setlength{\parskip}{6pt plus 2pt minus 1pt}}
}{% if KOMA class
  \KOMAoptions{parskip=half}}
\makeatother
\usepackage{color}
\usepackage{fancyvrb}
\newcommand{\VerbBar}{|}
\newcommand{\VERB}{\Verb[commandchars=\\\{\}]}
\DefineVerbatimEnvironment{Highlighting}{Verbatim}{commandchars=\\\{\}}
% Add ',fontsize=\small' for more characters per line
\usepackage{framed}
\definecolor{shadecolor}{RGB}{248,248,248}
\newenvironment{Shaded}{\begin{snugshade}}{\end{snugshade}}
\newcommand{\AlertTok}[1]{\textcolor[rgb]{0.94,0.16,0.16}{#1}}
\newcommand{\AnnotationTok}[1]{\textcolor[rgb]{0.56,0.35,0.01}{\textbf{\textit{#1}}}}
\newcommand{\AttributeTok}[1]{\textcolor[rgb]{0.13,0.29,0.53}{#1}}
\newcommand{\BaseNTok}[1]{\textcolor[rgb]{0.00,0.00,0.81}{#1}}
\newcommand{\BuiltInTok}[1]{#1}
\newcommand{\CharTok}[1]{\textcolor[rgb]{0.31,0.60,0.02}{#1}}
\newcommand{\CommentTok}[1]{\textcolor[rgb]{0.56,0.35,0.01}{\textit{#1}}}
\newcommand{\CommentVarTok}[1]{\textcolor[rgb]{0.56,0.35,0.01}{\textbf{\textit{#1}}}}
\newcommand{\ConstantTok}[1]{\textcolor[rgb]{0.56,0.35,0.01}{#1}}
\newcommand{\ControlFlowTok}[1]{\textcolor[rgb]{0.13,0.29,0.53}{\textbf{#1}}}
\newcommand{\DataTypeTok}[1]{\textcolor[rgb]{0.13,0.29,0.53}{#1}}
\newcommand{\DecValTok}[1]{\textcolor[rgb]{0.00,0.00,0.81}{#1}}
\newcommand{\DocumentationTok}[1]{\textcolor[rgb]{0.56,0.35,0.01}{\textbf{\textit{#1}}}}
\newcommand{\ErrorTok}[1]{\textcolor[rgb]{0.64,0.00,0.00}{\textbf{#1}}}
\newcommand{\ExtensionTok}[1]{#1}
\newcommand{\FloatTok}[1]{\textcolor[rgb]{0.00,0.00,0.81}{#1}}
\newcommand{\FunctionTok}[1]{\textcolor[rgb]{0.13,0.29,0.53}{\textbf{#1}}}
\newcommand{\ImportTok}[1]{#1}
\newcommand{\InformationTok}[1]{\textcolor[rgb]{0.56,0.35,0.01}{\textbf{\textit{#1}}}}
\newcommand{\KeywordTok}[1]{\textcolor[rgb]{0.13,0.29,0.53}{\textbf{#1}}}
\newcommand{\NormalTok}[1]{#1}
\newcommand{\OperatorTok}[1]{\textcolor[rgb]{0.81,0.36,0.00}{\textbf{#1}}}
\newcommand{\OtherTok}[1]{\textcolor[rgb]{0.56,0.35,0.01}{#1}}
\newcommand{\PreprocessorTok}[1]{\textcolor[rgb]{0.56,0.35,0.01}{\textit{#1}}}
\newcommand{\RegionMarkerTok}[1]{#1}
\newcommand{\SpecialCharTok}[1]{\textcolor[rgb]{0.81,0.36,0.00}{\textbf{#1}}}
\newcommand{\SpecialStringTok}[1]{\textcolor[rgb]{0.31,0.60,0.02}{#1}}
\newcommand{\StringTok}[1]{\textcolor[rgb]{0.31,0.60,0.02}{#1}}
\newcommand{\VariableTok}[1]{\textcolor[rgb]{0.00,0.00,0.00}{#1}}
\newcommand{\VerbatimStringTok}[1]{\textcolor[rgb]{0.31,0.60,0.02}{#1}}
\newcommand{\WarningTok}[1]{\textcolor[rgb]{0.56,0.35,0.01}{\textbf{\textit{#1}}}}
\usepackage{graphicx}
\makeatletter
\newsavebox\pandoc@box
\newcommand*\pandocbounded[1]{% scales image to fit in text height/width
  \sbox\pandoc@box{#1}%
  \Gscale@div\@tempa{\textheight}{\dimexpr\ht\pandoc@box+\dp\pandoc@box\relax}%
  \Gscale@div\@tempb{\linewidth}{\wd\pandoc@box}%
  \ifdim\@tempb\p@<\@tempa\p@\let\@tempa\@tempb\fi% select the smaller of both
  \ifdim\@tempa\p@<\p@\scalebox{\@tempa}{\usebox\pandoc@box}%
  \else\usebox{\pandoc@box}%
  \fi%
}
% Set default figure placement to htbp
\def\fps@figure{htbp}
\makeatother
\setlength{\emergencystretch}{3em} % prevent overfull lines
\providecommand{\tightlist}{%
  \setlength{\itemsep}{0pt}\setlength{\parskip}{0pt}}
\usepackage{bookmark}
\IfFileExists{xurl.sty}{\usepackage{xurl}}{} % add URL line breaks if available
\urlstyle{same}
\hypersetup{
  pdftitle={Docker and renv strategy},
  pdfauthor={R.G. Thomas},
  colorlinks=true,
  linkcolor={Maroon},
  filecolor={Maroon},
  citecolor={Blue},
  urlcolor={blue},
  pdfcreator={LaTeX via pandoc}}

\title{Docker and renv strategy}
\author{R.G. Thomas}
\date{}

\begin{document}
\maketitle

{
\hypersetup{linkcolor=}
\setcounter{tocdepth}{2}
\tableofcontents
}
\section{Introduction}\label{introduction}

Reproducibility is a cornerstone of professional data analysis, yet
achieving it with R Markdown can be challenging. R projects often break
on other computers due to mismatched R versions or package versions,
sending developers into ``dependency hell.'' To solve this, we can
leverage \textbf{renv} (for R package management) and \textbf{Docker}
(for containerizing the computing environment). Together, these tools
ensure that an R Markdown document runs anywhere with the same packages,
R version, and system libraries as the original setup. Below, we present
a structured guide on using renv and Docker for a reproducible R
Markdown workflow, with diagrams, code snippets, and clear steps.

R is an open-source programming language for statistics and data
science. The \textbf{renv} package helps create reproducible R project
environments by isolating package dependencies.

\begin{center}\rule{0.5\linewidth}{0.5pt}\end{center}

\subsection{Using renv for R Package
Reproducibility}\label{using-renv-for-r-package-reproducibility}

\textbf{renv} (Reproducible Environment) is an R package that manages
library dependencies for your project. Instead of sharing one system
library across projects, renv gives each project its own library of R
packages and a lockfile (\texttt{renv.lock}) recording exact package
versions. This means that if two developers use renv, they can install
identical package versions for a project, avoiding version conflicts and
``it works on my machine'' issues. In practice, renv makes it easy to
\textbf{snapshot} your R packages and later \textbf{restore} them on any
system.

\subsubsection{Key Features of renv:}\label{key-features-of-renv}

\begin{itemize}
\tightlist
\item
  \textbf{Isolated project library}: renv creates a project-specific
  library (usually in \texttt{renv/library}) containing all packages
  used by that project. This isolation means updates in one project
  won't affect others.
\item
  \textbf{Lockfile for dependencies}: When you finish installing or
  updating packages, you call \texttt{renv::snapshot()}. This produces
  \textbf{\texttt{renv.lock}}, a JSON file listing each package and the
  exact version (and source) in use. The lockfile is meant to be
  committed to version control and shared.
\item
  \textbf{Restoring environment}: On a new machine (or when reproducing
  past results), you use \texttt{renv::restore()} to install the
  \textbf{exact versions} of packages from the lockfile. This gives an R
  environment \textbf{identical} to the one that created the lockfile,
  as long as the same R version is available.
\end{itemize}

\subsubsection{Example: Using renv in an R
Project}\label{example-using-renv-in-an-r-project}

\begin{Shaded}
\begin{Highlighting}[]
\CommentTok{\# Install and initialize renv in a new R project}
\FunctionTok{install.packages}\NormalTok{(}\StringTok{"renv"}\NormalTok{)      }\CommentTok{\# one{-}time installation of renv}
\NormalTok{renv}\SpecialCharTok{::}\FunctionTok{init}\NormalTok{()                  }\CommentTok{\# initialize renv (creates renv infrastructure)}

\CommentTok{\# ...After installing or updating some project{-}specific packages...}
\NormalTok{renv}\SpecialCharTok{::}\FunctionTok{snapshot}\NormalTok{()              }\CommentTok{\# save exact package versions to renv.lock}

\CommentTok{\# ...Later or on another system...}
\NormalTok{renv}\SpecialCharTok{::}\FunctionTok{restore}\NormalTok{()               }\CommentTok{\# restore packages from renv.lock for reproducibility}
\end{Highlighting}
\end{Shaded}

In an R Markdown workflow, you would: attach renv to the project, work
on your \texttt{.Rmd} using whatever CRAN/Bioconductor packages needed,
and finally snapshot the environment. The \texttt{renv.lock} file
ensures that anyone with the lockfile can recreate the \textbf{same R
package environment}. However, renv alone does \textbf{not lock the R
version or system libraries} -- that's where Docker comes in.

\begin{center}\rule{0.5\linewidth}{0.5pt}\end{center}

\subsection{Using Docker for OS-Level
Reproducibility}\label{using-docker-for-os-level-reproducibility}

While renv handles R packages, \textbf{Docker} takes care of the
\textbf{operating system, R interpreter, and system libraries}. A Docker
container is like a lightweight virtual machine image that includes
everything needed to run your project: the specific OS (Linux
distribution), exact R version, required system libraries (e.g.~C/C++
libraries), and your R code. By running an R Markdown project in Docker,
you ensure that differences in OS or R installation are no longer an
issue -- any machine running Docker will run the container
\emph{identically}.

Docker allows you to \textbf{build an image} that encapsulates your
environment. For R workflows, a common practice is to start from a base
image (for example, the Rocker project provides images like
\texttt{rocker/r-ver:\textless{}R-version\textgreater{}} for R). You
then add your content and dependencies.

\subsubsection{Example: A Simple Dockerfile for an R Markdown
Project}\label{example-a-simple-dockerfile-for-an-r-markdown-project}

\begin{Shaded}
\begin{Highlighting}[]
\CommentTok{\# Use R 4.1.0 on Linux as base image}
\KeywordTok{FROM}\NormalTok{ rocker/r{-}ver:4.1.0  }

\CommentTok{\# Copy renv lockfile into the image}
\KeywordTok{COPY}\NormalTok{ renv.lock /renv.lock}

\CommentTok{\# Install renv and restore packages from lockfile}
\KeywordTok{RUN} \ExtensionTok{R} \AttributeTok{{-}e} \StringTok{"install.packages(\textquotesingle{}renv\textquotesingle{}, repos=\textquotesingle{}https://cloud.r{-}project.org\textquotesingle{}); renv::restore()"}  

\CommentTok{\# Copy the rest of the project (e.g., R Markdown files)}
\KeywordTok{COPY}\NormalTok{ . /workspace}

\CommentTok{\# Set working directory and default command (if needed)}
\KeywordTok{WORKDIR}\NormalTok{ /workspace}
\KeywordTok{CMD}\NormalTok{ [}\StringTok{"Rscript"}\NormalTok{, }\StringTok{"{-}e"}\NormalTok{, }\StringTok{"rmarkdown::render(\textquotesingle{}analysis.Rmd\textquotesingle{})"}\NormalTok{]}
\end{Highlighting}
\end{Shaded}

In the snippet above, the Docker image will install the exact package
versions from \texttt{renv.lock}. The \texttt{renv::restore()} command
inside \texttt{RUN} will install all the R packages (and exact versions)
listed in the lockfile into the image's R library.

\begin{center}\rule{0.5\linewidth}{0.5pt}\end{center}

\subsection{Combining renv and Docker in R Markdown
Workflows}\label{combining-renv-and-docker-in-r-markdown-workflows}

Using renv or Docker alone improves reproducibility, but
\textbf{combining them gives the best of both worlds}. Docker will
guarantee the OS and R version, and renv will guarantee the R packages.
Together, you achieve \textbf{end-to-end reproducibility} --- from
operating system all the way to the R code's package versions.

\subsubsection{Steps to Integrate renv with
Docker:}\label{steps-to-integrate-renv-with-docker}

\begin{enumerate}
\def\labelenumi{\arabic{enumi}.}
\tightlist
\item
  \textbf{Develop with renv}: Start a new R project for your R Markdown
  analysis and run \texttt{renv::init()}. Install all the R packages
  your \texttt{.Rmd} needs. Periodically run \texttt{renv::snapshot()}
  to update the lockfile.
\item
  \textbf{Write a Dockerfile}: Define the base image (e.g.,
  \texttt{FROM\ rocker/r-ver:4.1.0}), copy the \texttt{renv.lock},
  install dependencies, and set up the project.
\item
  \textbf{Build the Docker image}: Run
  \texttt{docker\ build\ -t\ myanalysis/image:v1\ .} to create the
  image.
\item
  \textbf{Share the image}: Push it to Docker Hub or share the file, so
  others can pull and run it.
\item
  \textbf{Run anywhere with confidence}: The container ensures
  consistency across machines without additional manual setup.
\end{enumerate}

\begin{center}\rule{0.5\linewidth}{0.5pt}\end{center}

\subsection{Conclusion}\label{conclusion}

By incorporating \textbf{renv} and \textbf{Docker}, an R Markdown
project becomes \textbf{easy to share and fully reproducible}. The
original author can provide a Docker image (or Dockerfile) along with
the R Markdown and \texttt{renv.lock}. A collaborator or reviewer can
launch the container and get the \textbf{same results}, without worrying
about package versions or system setup. This approach is considered a
\textbf{best practice for long-term reproducibility in R}.

Together, \textbf{renv and Docker} ensure that your R Markdown analyses
are \textbf{completely reproducible and portable}, meeting the high
standards required for professional data science and research
documentation.

\end{document}
